% EGA V, English translation by P. Blass, 1987 (cote 24): archive cote 24, site page 17 (PDF page 18).
% typescript. Draft compared: none.
% AI-transcribed by Claude Opus 5.5 (claude-opus-5-5, adaptive thinking, effort high), Message Batches API, batch ega24b (msgbatch_01RGJJeHdR96U16ZWFuUjtDx), 2026-11-13; not yet checked by eye.
\scanpage{24/17}
\mspage{17}

then $\underline{P}(O_{X,x}, k) \Leftrightarrow \underline{P}(O_{Y_\eta,y}, k(y))$. Indeed taking into account par.
6* (iii) and (iv) follow from (ii) and 2.8 (ii). On the other
hand (i) has been proven in 2.12 and it remains to deduce (ii)
from (i). To do this let Z be the set of points where X is not
smooth over k, its inverse image $Z_\eta$ in $Y_\eta$ is therefore (by 2.12)
the set of points of $Y_\eta$ at which $Y_\eta$ is not smooth over $K(\eta)$. But
the codimension of Z in X at x is equal to that of $Z_\eta$ in $Y_\eta$ at y
because of flatness (reference par. 6*). Therefore one is $>k$ if
the other one is such which completes the proof.

\uline{Corollary 2.15}

Let Z be the set of points of X at which X is not smooth over
k (respectively is not geometrically $R_k$ over k, respectively is not
separable over k, respectively is not normal over k). In order
for $Y_\eta$ to be smooth (respectively geometrically $R_k$, respectively
separable, respectively geometrically normal) over $k(\eta)$ it is
necessary and sufficient that Z should be finite.

From writing up point of view statements 2.14 and 2.15
contain 2.12 and 2.13 (which we could thus dispense with stating
separately) on the other hand the corollary is practically more
important than the theorem which one could give in a proposition
or a preliminary lemma so that the corollary would be more
glorified. We can modify the case of property (iii):

\uline{Corollary 2.16}

(Let us still suppose that f is an immersion and also that F
is coherent), with the conditions of 2.7, in order that $Z_i$ should

\margin{* Ask Grothendieck}
